\documentclass[11pt]{article}
\usepackage[margin=1in]{geometry}
\usepackage{amsmath,amssymb}
\usepackage{booktabs,longtable,array}
\usepackage{enumitem}
\usepackage{xcolor}
\usepackage{listings}
\usepackage{hyperref}
\usepackage{microtype}
\hypersetup{colorlinks=true,linkcolor=black,urlcolor=blue}
\lstset{basicstyle=\ttfamily\small,frame=single,columns=fullflexible,keepspaces=true,showstringspaces=false,breaklines=true}
\newcommand{\codepath}[1]{\nolinkurl{#1}}
\newcommand{\R}{\mathbb{R}}
\title{\textbf{AutoIA Software Manual}\\
User Guide, Operations Reference, and Developer Maintenance Guide\\
Version 0.21.0.dev1}
\author{}
\date{August 2026}
\begin{document}
\maketitle

\begin{abstract}
This manual explains how to install, use, diagnose, test, and modify AutoIA
version 0.21.0.dev1.  It is written for two overlapping audiences.  The first is
an engineering user who wants to run validated interval analysis on ordinary
Python/NumPy code.  The second is a researcher/developer who needs to understand
the source layout well enough to add a function, representation, semantic
primitive, controller policy, performance backend, or many-box execution feature
without violating the certification architecture.

All code examples are complete Python fragments intended to be directly
copyable.  The manual states the supported-language boundary explicitly; code
that happens to work outside that boundary should not be treated as a stable
Version-1 contract.
\end{abstract}

\tableofcontents
\newpage

%===============================================================================
\section{Installation and project layout}
%===============================================================================

AutoIA 0.21.0.dev1 requires Python 3.10 or newer and NumPy 1.24 or newer.  SciPy
is optional; it is needed only if the user's model itself imports SciPy and wants
the supported SciPy normal/error-function syntax.

To install the built wheel into an environment:
\begin{lstlisting}[language=bash]
python -m pip install autoia-0.21.0.dev1-py3-none-any.whl
\end{lstlisting}

To verify the installed version:
\begin{lstlisting}[language=Python]
import autoia
print(autoia.__version__)
# 0.21.0.dev1
\end{lstlisting}

The most important public imports are:
\begin{lstlisting}[language=Python]
from autoia import (
    autoia,
    autoia_subdivide,
    autoia_adaptive_subdivide,
    prepare,
    Interval,
    RigorousBinary64Backend,
    RangePolicy,
    SubdivisionPolicy,
    BranchPolicy,
)
\end{lstlisting}

%===============================================================================
\section{First validated calculation}
%===============================================================================

The smallest useful example is:
\begin{lstlisting}[language=Python]
import numpy as np
from autoia import autoia, Interval


def model(p):
    x = p[0]
    return np.exp(x) + x*x


result = autoia(
    model,
    [Interval(-0.2, 0.3)],
    order=4,
)

print("lower =", result.lower)
print("upper =", result.upper)
print("width =", result.width)
print("machine rigorous =", result.machine_rigorous)
\end{lstlisting}

If \texttt{autoia(...)} returns normally, the reported interval is intended to
contain the exact real program output for every input in the supplied interval.
If the requested statement cannot be certified, the function raises a typed
\texttt{AutoIAError} subclass and returns no rigorous result.

%===============================================================================
\section{Specifying uncertain inputs}
%===============================================================================

For one uncertain scalar parameter use
\begin{lstlisting}[language=Python]
Interval(lower, upper)
\end{lstlisting}
with finite endpoints satisfying \texttt{lower <= upper}.

For \(m\) uncertain inputs, pass a sequence of \(m\) intervals:
\begin{lstlisting}[language=Python]
box = [
    Interval(90.0, 110.0),
    Interval(900.0, 1100.0),
    Interval(9.5, 10.5),
]
result = autoia(model, box)
\end{lstlisting}

The uncertainty model is currently a Cartesian box.  Correlated physical input
sets such as ellipsoids, simplices encoded intrinsically, or arbitrary
non-rectangular feasible sets are not yet first-class input-domain objects.  A
user may still write algebraic relations in the model, but the declared uncertain
parameters themselves are supplied as independent interval coordinates.

%===============================================================================
\section{Choosing a representation}
%===============================================================================

The public \texttt{representation=} option accepts:

\begin{longtable}{p{0.22\textwidth}p{0.68\textwidth}}
\toprule
Value & Meaning and recommended use \\
\midrule
\endhead
\texttt{"polynomial"} & Default.  Retains a sparse polynomial in the original uncertain inputs plus a rigorous interval remainder.  Usually the best starting point for smooth engineering models. \\
\texttt{"generator"} & Retains the same physical-input polynomial plus shared auxiliary error generators.  Use when downstream reuse/cancellation of approximated vector/intermediate quantities matters. \\
\texttt{"interval"} & Uses ordinary interval enclosures without retaining explicit polynomial dependence.  Robust fallback and useful near some exact function-domain boundaries, but usually less sharp for repeated dependence. \\
\bottomrule
\end{longtable}

Example:
\begin{lstlisting}[language=Python]
result = autoia(
    model,
    [Interval(-1.0, 1.0)],
    order=3,
    representation="generator",
)

print(result.generator_ids)
print(result.generator_matrix)
\end{lstlisting}

The generator representation is not automatically superior in every model.  It
preserves additional error provenance at the cost of extra bookkeeping and can
accumulate many generator IDs in long nonlinear computations.

%===============================================================================
\section{Polynomial order}
%===============================================================================

The \texttt{order} argument is a nonnegative integer controlling the maximum
total degree retained in the polynomial representation.

\begin{lstlisting}[language=Python]
result = autoia(model, box, order=4)
\end{lstlisting}

Higher order can reduce the interval remainder by retaining more physical-input
dependence, but increases polynomial coefficient count and range-analysis cost.
The appropriate value is model-dependent.  Orders around 2--5 are the natural
starting range for many small engineering examples in the current prototype.

Changing \texttt{order} affects efficiency and sharpness, not the mathematical
inclusion contract: discarded polynomial terms are rigorously transferred to
remainder bounds.

%===============================================================================
\section{Supported scalar arithmetic}
%===============================================================================

The core supported scalar operations include:
\begin{itemize}[leftmargin=*]
\item addition and subtraction;
\item multiplication;
\item division when the denominator can be proved separated from zero;
\item integer powers, including negative integer powers when zero is excluded;
\item constant non-integer real powers on the supported positive-real branch.
\end{itemize}

Examples:
\begin{lstlisting}[language=Python]
def model(p):
    x = p[0]
    a = x*x + 2.0*x - 1.0
    b = 1.0 / (2.0 + x*x)
    c = (1.5 + x)**1.25
    d = (2.0 + x)**-0.5
    return np.array([a, b, c, d])
\end{lstlisting}

An uncertain exponent such as \texttt{x ** p[1]} is currently unsupported.

%===============================================================================
\section{Supported elementary and statistical functions}
%===============================================================================

The ordinary NumPy and standard-library forms of these scalar functions are
supported when visible through the model's module/global namespace:
\begin{lstlisting}[language=Python]
np.sqrt(x)
np.exp(x)
np.log(x)
np.sin(x)
np.cos(x)

math.sqrt(x)
math.exp(x)
math.log(x)
math.sin(x)
math.cos(x)
math.erf(x)
\end{lstlisting}

AutoIA also exports:
\begin{lstlisting}[language=Python]
from autoia import (
    erf,
    standard_normal_pdf,
    standard_normal_cdf,
    normal_pdf,
    normal_cdf,
    norm_pdf,
    norm_cdf,
)
\end{lstlisting}

For a normal distribution with constant mean \texttt{mu} and standard deviation
\texttt{sigma > 0}:
\begin{lstlisting}[language=Python]
def model(p):
    x = p[0]
    return np.array([
        normal_pdf(x, mean=10.0, std=2.0),
        normal_cdf(x, mean=10.0, std=2.0),
    ])
\end{lstlisting}

The current normal convenience APIs require constant distribution parameters.
Uncertain mean/standard deviation parameters are not yet supported by those
front-end helpers.

%===============================================================================
\section{SciPy compatibility}
%===============================================================================

If SciPy is installed and the user model imports it, these common forms are
recognized:
\begin{lstlisting}[language=Python]
import scipy

scipy.special.erf(x)
scipy.special.ndtr(x)

scipy.stats.norm.pdf(x)
scipy.stats.norm.cdf(x)
scipy.stats.norm.pdf(x, loc=mu, scale=sigma)
scipy.stats.norm.cdf(x, loc=mu, scale=sigma)
\end{lstlisting}

Frozen legacy normal distributions are supported:
\begin{lstlisting}[language=Python]
from scipy.stats import norm
rv = norm(loc=2.0, scale=1.5)


def model(p):
    x = p[0]
    return np.array([rv.pdf(x), rv.cdf(x)])
\end{lstlisting}

When available in the installed SciPy version, the newer
\texttt{scipy.stats.Normal} object API is also recognized for PDF/CDF use.

SciPy does not become proof code.  On validated inputs the local dispatcher
routes these calls to AutoIA's own rigorous error-function/normal kernels.
Numeric-only calls continue to use ordinary SciPy.

%===============================================================================
\section{NumPy arrays and ordinary engineering assembly}
%===============================================================================

Inside AutoIA's private execution namespace, common array constructors default to
\texttt{dtype=object} when the user did not explicitly request a dtype.  This
allows standard engineering assembly:
\begin{lstlisting}[language=Python]
import numpy as np


def model(p):
    k1, k2, force = p

    A = np.zeros((2, 2))
    b = np.zeros(2)

    A[0, 0] = k1 + k2
    A[0, 1] = -k2
    A[1, 0] = -k2
    A[1, 1] = k2

    b[0] = force

    return np.linalg.solve(A, b)
\end{lstlisting}

Supported constructors include common forms of \texttt{array}, \texttt{asarray},
\texttt{zeros}, \texttt{ones}, \texttt{empty}, \texttt{full}, their
\texttt{*\_like} variants, \texttt{eye}, and \texttt{identity}.

Ordinary object-array indexing, slicing, reshape, transpose, stacking,
concatenation, and item assignment are supported in the tested engineering cases.

%===============================================================================
\section{Broadcasting, reductions, and matrix composition}
%===============================================================================

Supported broadcasted scalar arithmetic includes cases where a numeric constant
NumPy array created outside the model is combined with one validated scalar:
\begin{lstlisting}[language=Python]
COEFF = np.array([1.5, -2.0, 0.25])


def model(p):
    return COEFF * p[0]
\end{lstlisting}

The supported reduction subset includes:
\begin{lstlisting}[language=Python]
np.sum(A, axis=...)
np.prod(A, axis=...)
np.mean(A, axis=...)
np.trace(A)
\end{lstlisting}
with common \texttt{axis} and \texttt{keepdims} behavior.  For validated
reductions, omit \texttt{dtype} or use \texttt{dtype=object}.  General
\texttt{where=} and \texttt{out=} semantics are not frozen.

The tested 1D/2D composition subset includes:
\begin{lstlisting}[language=Python]
np.dot(a, b)
np.matmul(A, x)
A @ x
np.inner(a, b)
np.outer(a, b)
\end{lstlisting}

%===============================================================================
\section{Norms}
%===============================================================================

The supported norm forms are:
\begin{lstlisting}[language=Python]
np.linalg.norm(x)             # vector Euclidean norm
np.linalg.norm(x, ord=2)      # vector Euclidean norm
np.linalg.norm(A)             # matrix Frobenius norm
np.linalg.norm(A, ord="fro") # matrix Frobenius norm
\end{lstlisting}

For a vector \(x=(x_1,\ldots,x_n)^T\), the supported Euclidean norm is
\[
\|x\|_2=\sqrt{\sum_{i=1}^n x_i^2}.
\]
Matrix spectral \texttt{ord=2}, general \(p\)-norms, nuclear norms, and max-based
norms are not currently supported.

%===============================================================================
\section{Semantic linear solves}
%===============================================================================

Use ordinary NumPy syntax:
\begin{lstlisting}[language=Python]
u = np.linalg.solve(A, b)
\end{lstlisting}

AutoIA recognizes this call and applies its semantic linear-system certificate.
The entries of \texttt{A} and \texttt{b} may contain retained polynomial
nonlinearity from earlier arithmetic or elementary functions.

Multiple right-hand sides are supported:
\begin{lstlisting}[language=Python]
X = np.linalg.solve(A, B)
\end{lstlisting}
where \texttt{B} is a matrix with one right-hand side per column.

The solve may fail with \texttt{CertificationError} if AutoIA cannot prove the
matrix family nonsingular/contractive on the supplied box.  This is not
necessarily evidence that the true matrix is singular.  Subdivision, a smaller
input box, or a later stronger certifier may succeed.

%===============================================================================
\section{Uncertain Python branches}
%===============================================================================

Ordinary scalar conditionals are supported for the documented subset:
\begin{lstlisting}[language=Python]
def model(p):
    x = p[0]

    if x < -0.25:
        y = -1.0 + x*x
    elif x <= 0.5:
        y = x*x
    else:
        y = 1.0 + x

    return y*y + 0.1*x
\end{lstlisting}

Supported scalar comparisons include \texttt{<}, \texttt{<=}, \texttt{>},
\texttt{>=}, \texttt{==}, and \texttt{!=}, plus scalar truth testing
\texttt{if x:}.  Python short-circuit \texttt{and}/\texttt{or} works through the
same scalar truth mechanism in the tested cases.

If a predicate cannot be rigorously decided on the current box, AutoIA may split
the input box or enumerate both feasible execution paths.  A branching model
returns an \texttt{AutoIAPathwiseResult} with branch metadata in addition to
global lower/upper bounds.

%===============================================================================
\section{Branch policy and resource limits}
%===============================================================================

The default branch policy is equivalent to:
\begin{lstlisting}[language=Python]
from autoia import BranchPolicy

policy = BranchPolicy(
    max_branches=64,
    max_path_decisions=24,
    max_split_depth_per_decision=2,
    max_total_splits=64,
    prefer_geometric_split=True,
)
\end{lstlisting}

Use it with:
\begin{lstlisting}[language=Python]
result = autoia(model, box, branch_policy=policy)
\end{lstlisting}

If a hard path resource budget is exhausted before every feasible path is
certified, AutoIA raises \texttt{ResourceLimitError}.  It does not silently drop
paths.

%===============================================================================
\section{Range policies}
%===============================================================================

The default final range policy is adaptive.  The public policy object is:
\begin{lstlisting}[language=Python]
from autoia import RangePolicy

policy = RangePolicy(
    internal_method="termwise",
    final_method="adaptive",
    bernstein_max_coefficients=256,
    proof_bernstein_max_coefficients=256,
)
\end{lstlisting}

The supported internal methods are:
\begin{itemize}[leftmargin=*]
\item \texttt{"termwise"};
\item \texttt{"monotonicity"};
\item \texttt{"bernstein"};
\item \texttt{"portfolio"}.
\end{itemize}
The final method may additionally be \texttt{"adaptive"}.

Use either a \texttt{RangePolicy} object or the legacy convenience keyword
options such as \texttt{range\_method}; do not pass both styles in the same call.

%===============================================================================
\section{Prescribed subdivision}
%===============================================================================

To force a Cartesian subdivision, use:
\begin{lstlisting}[language=Python]
from autoia import autoia_subdivide

piecewise = autoia_subdivide(
    model,
    [Interval(-1.0, 1.0), Interval(0.0, 2.0)],
    split_counts=[4, 2],
    order=4,
)
\end{lstlisting}

The \texttt{split\_counts} sequence contains one positive integer per input
coordinate.  Each resulting subbox is independently certified.

%===============================================================================
\section{Adaptive subdivision}
%===============================================================================

Use:
\begin{lstlisting}[language=Python]
from autoia import autoia_adaptive_subdivide, SubdivisionPolicy

policy = SubdivisionPolicy(
    max_branches=16,
    max_depth=6,
    max_probe_directions=2,
    min_structural_gain=0.05,
    min_relative_width_gain=1.0e-3,
    split_method="cheap",
)

result = autoia_adaptive_subdivide(
    model,
    box,
    order=4,
    subdivision_policy=policy,
)
\end{lstlisting}

The release-oriented default \texttt{split\_method="cheap"} scores hypothetical
splits from the existing polynomial and only certifies the selected child pair.
Use \texttt{split\_method="lookahead"} to reproduce the expensive historical
one-step lookahead research oracle.

The current adaptive-subdivision controller does not generally compose with
uncertain Python control flow.  For path-dependent models use ordinary path-aware
\texttt{autoia(...)} or prescribed \texttt{autoia\_subdivide(...)} unless the
current documentation explicitly states otherwise.

%===============================================================================
\section{Understanding result objects}
%===============================================================================

For an ordinary result:
\begin{lstlisting}[language=Python]
print(result.lower)
print(result.upper)
print(result.width)
print(result.output_shape)
print(result.lower_shaped)
print(result.upper_shaped)
print(result.width_shaped)
print(result.machine_rigorous)
print(result.representation_names)
\end{lstlisting}

The flat \texttt{lower}/\texttt{upper} arrays are retained for compatibility.
Use the shaped properties to recover the model's returned array shape.

In generator mode:
\begin{lstlisting}[language=Python]
print(result.generator_ids)
print(result.generator_matrix)
\end{lstlisting}
Rows of \texttt{generator\_matrix} correspond to flattened outputs and columns
correspond to shared generator IDs.

The \texttt{objects} field contains the validated scalar representations.  It is
intended primarily for research/diagnostic use; ordinary users usually need only
bounds and diagnostics.

%===============================================================================
\section{Understanding pathwise results}
%===============================================================================

A path-aware result additionally exposes information such as:
\begin{lstlisting}[language=Python]
print(result.branch_results)
print(result.branch_boxes)
print(result.branch_decisions)
print(result.branch_traces)
print(result.path_events)
\end{lstlisting}

These records are useful when a wide global interval is caused by genuinely
different program regimes.  They should also be examined when path budgets become
large or a branch-local singularity prevents certification.

%===============================================================================
\section{Prepared sessions for repeated queries}
%===============================================================================

For repeated certification of the same program/configuration, use \texttt{prepare}
instead of calling one-shot \texttt{autoia} repeatedly:
\begin{lstlisting}[language=Python]
from autoia import prepare, Interval

session = prepare(
    model,
    order=4,
    representation="generator",
    cache_max_entries=256,
)

r1 = session.certify([
    Interval(-1.0, 1.0),
    Interval(0.0, 2.0),
])

r2 = session.certify([
    Interval(-0.25, 0.25),
    Interval(0.5, 1.0),
])

print(session.stats)
print(r2.cache_info)
\end{lstlisting}

The session prepares the model's local dispatch once and reuses one audited
trusted backend.

%===============================================================================
\section{Certificate cache behavior}
%===============================================================================

The session cache has three important query outcomes.

\paragraph{Fresh certificate.}
\begin{lstlisting}[language=Python]
result.cache_info["kind"] == "fresh-certificate"
\end{lstlisting}
The engineering model was executed and a new proof was produced.

\paragraph{Exact certificate hit.}
\begin{lstlisting}[language=Python]
result.cache_info["kind"] == "exact-certificate"
\end{lstlisting}
The exact same stored input endpoints were already certified; the model is not
re-executed.

\paragraph{Restricted certificate hit.}
\begin{lstlisting}[language=Python]
result.cache_info["kind"] == "restricted-certificate"
\end{lstlisting}
A previously certified parent box contained the new query, and AutoIA rigorously
re-centered/re-ranged the stored polynomial/generator certificate on the new
interior box without executing the model.

The cache is bounded by \texttt{cache\_max\_entries}.  Set it to zero to disable
persistent cache storage:
\begin{lstlisting}[language=Python]
session = prepare(model, cache_max_entries=0)
\end{lstlisting}

Use:
\begin{lstlisting}[language=Python]
session.clear_cache()
\end{lstlisting}
to remove stored certificates.

%===============================================================================
\section{Many-box finite collections}
%===============================================================================

A finite collection may be supplied as nested endpoint arrays:
\begin{lstlisting}[language=Python]
boxes = np.array([
    [[-1.0, -0.8], [0.0, 0.2]],
    [[-0.5, -0.3], [0.4, 0.6]],
    [[ 0.1,  0.2], [0.8, 1.0]],
], dtype=float)

results = session.certify_many(
    boxes,
    parallel="serial",
)
\end{lstlisting}

or as ordinary lists of \texttt{Interval} objects.

For bounds-only output:
\begin{lstlisting}[language=Python]
lower, upper = session.certify_many_bounds(boxes)
print(lower.shape)
print(upper.shape)
\end{lstlisting}
If there are \(N\) boxes and \(r\) flattened scalar outputs, the arrays have
shape \((N,r)\).

%===============================================================================
\section{CPU parallel many-box execution}
%===============================================================================

The supported modes are:
\begin{lstlisting}[language=Python]
parallel="serial"
parallel="thread"
parallel="process"
\end{lstlisting}

Example process mode:
\begin{lstlisting}[language=Python]
results = session.certify_many(
    boxes,
    parallel="process",
    workers=8,
    chunk_size=64,
    use_cache=True,
    allow_restriction=True,
)
\end{lstlisting}

The model must be pickleable for process mode.  The safest pattern is a
module-level function:
\begin{lstlisting}[language=Python]
# good for multiprocessing

def engineering_model(p):
    ...
\end{lstlisting}
A nested closure or lambda is generally not pickleable under the spawn execution
mode.  Use serial/thread mode or move the model to module scope.

Each worker process prepares the model once and creates its own audited trusted
backend.  Session-local caches are not automatically shared across processes.
For heavily overlapping boxes, one serial/shared session with certificate reuse
may outperform process mode despite using fewer cores; benchmark the actual
workload.

%===============================================================================
\section{Streaming millions of boxes}
%===============================================================================

Do not construct a list of one million full result objects if the downstream
workflow only needs bounds.  Use a generator and the streaming API:
\begin{lstlisting}[language=Python]
from autoia import Interval, prepare


def generate_boxes():
    for i in range(1_000_000):
        a = -1.0 + 2.0*i/1_000_000
        yield [Interval(a, a + 1.0e-3)]


session = prepare(model, order=3)

for lower, upper in session.iter_bounds_many(
    generate_boxes(),
    parallel="process",
    workers=8,
    chunk_size=256,
):
    # Write to disk/database or update an aggregate.
    consume(lower, upper)
\end{lstlisting}

The iterator keeps one executor alive and only a bounded number of chunks in
flight.  This controls orchestration memory, although each individual AutoIA
certification still uses the current Python-object representation internally.

%===============================================================================
\section{Session statistics}
%===============================================================================

Inspect:
\begin{lstlisting}[language=Python]
print(session.stats)
\end{lstlisting}
The returned \texttt{PreparedSessionStats} fields are:
\begin{itemize}[leftmargin=*]
\item \texttt{model\_preparations};
\item \texttt{certify\_requests};
\item \texttt{engineering\_model\_executions};
\item \texttt{exact\_cache\_hits};
\item \texttt{restricted\_cache\_hits};
\item \texttt{cache\_misses};
\item \texttt{cache\_entries}.
\end{itemize}
These statistics are especially useful for deciding whether a repeated-query
workload is benefiting from mathematical certificate reuse or merely from CPU
parallelism.

%===============================================================================
\section{Trusted reduction backend configuration}
%===============================================================================

The default backend chooses the historical scalar trusted reduction for short
reductions and the vectorized outward pairwise implementation for longer ones.
You can create an explicit backend:
\begin{lstlisting}[language=Python]
from autoia import RigorousBinary64Backend

backend = RigorousBinary64Backend(
    reduction_method="auto",
    vectorized_reduction_min_terms=16,
)

result = autoia(model, box, backend=backend)
\end{lstlisting}

For research comparisons:
\begin{lstlisting}[language=Python]
reference_backend = RigorousBinary64Backend(
    reduction_method="reference"
)

fast_backend = RigorousBinary64Backend(
    reduction_method="vectorized"
)
\end{lstlisting}

Different rigorous reduction orders can produce slightly different certified
endpoints.  Bitwise identity is not a general requirement; containment is.

%===============================================================================
\section{Typed errors and how to respond}
%===============================================================================

Catch the common base class if an application wants one failure handler:
\begin{lstlisting}[language=Python]
from autoia import AutoIAError

try:
    result = autoia(model, box)
except AutoIAError as exc:
    print(type(exc).__name__)
    print(exc)
\end{lstlisting}

Use the subtype to diagnose the problem:

\begin{longtable}{p{0.28\textwidth}p{0.62\textwidth}}
\toprule
Exception & Meaning / likely response \\
\midrule
\endhead
\texttt{InvalidInputError} & Malformed public call, invalid interval/order/policy/backend option.  Correct the call. \\
\texttt{UnsupportedOperationError} & Program used syntax/operation outside the current supported subset.  Rewrite using supported operations or implement a validated extension. \\
\texttt{CertificationError} & Supported mathematical operation was recognized but the requested proof failed, e.g. zero-crossing division or failed solve contraction.  Subdivide or improve the certifier. \\
\texttt{ResourceLimitError} & Path/subdivision budget exhausted.  Increase a justified budget or reduce the box. \\
\texttt{NumericalLimitError} & Trusted finite computation overflow/underflow or environment numerical limitation prevented a proof. \\
\texttt{ModelExecutionError} & The user's ordinary Python model raised an exception or returned incompatible path shapes. \\
\texttt{InternalAutoIAError} & AutoIA detected an internal inconsistency; treat this as a software defect and preserve a reproducible case. \\
\bottomrule
\end{longtable}

%===============================================================================
\section{Current supported-language boundaries}
%===============================================================================

Do not rely on the following as stable validated semantics unless later releases
explicitly add them:
\begin{itemize}[leftmargin=*]
\item \texttt{np.max}, \texttt{np.min}, \texttt{np.maximum}, \texttt{np.minimum};
\item general \texttt{abs}/sign/nonsmooth array operations;
\item uncertain Boolean masks and \texttt{np.where} with uncertain masks;
\item uncertain array indices or slices;
\item arbitrary \texttt{einsum}/\texttt{tensordot} signatures;
\item spectral/eigenvalue matrix norms;
\item general uncertainty-dependent \texttt{while} loop termination;
\item functions hidden inside compiled extensions or helpers whose globals are not
visible to the local shadow dispatcher;
\item arbitrary runtime-local imports inside the model body;
\item general nonlinear root/ODE/eigen/optimization semantic primitives.
\end{itemize}

The correct behavior for unsupported semantics is fail closed, not silently
coerce validated objects to ordinary floating point.

%===============================================================================
\section{Known boundary issue for sqrt/log}
%===============================================================================

A user interval may mathematically lie exactly on a valid domain boundary, such
as \([0,4]\) for \(\sqrt{x}\).  The polynomial representation uses a rigorous
symmetric centered envelope that can extend by a tiny amount below the original
lower endpoint because of outward rounding.  AutoIA will not silently clip that
negative proof interval.

If this occurs, practical options are:
\begin{itemize}[leftmargin=*]
\item use the interval representation for that calculation if appropriate;
\item subdivide or move the physical domain away from the exact boundary;
\item wait for the planned richer boundary/domain representation before relying
on polynomial-mode exact-boundary behavior.
\end{itemize}

%===============================================================================
\section{Performance guidance for users}
%===============================================================================

For one or a few boxes:
\begin{itemize}[leftmargin=*]
\item start with \texttt{representation="polynomial"};
\item use the default backend/range policy;
\item increase polynomial order only when remainder/range width justifies the
extra cost;
\item inspect branching/subdivision only when the global result is too wide or a
proof fails.
\end{itemize}

For many boxes over one program:
\begin{itemize}[leftmargin=*]
\item use \texttt{prepare};
\item keep overlapping/interior queries in one session so certificate restriction
can occur;
\item inspect \texttt{session.stats};
\item use process mode for many unrelated fresh boxes on a multicore CPU;
\item use streaming bounds APIs for very large collections;
\item benchmark cache reuse versus process parallelism on the actual workload.
\end{itemize}

The current implementation is not yet a true vectorized million-box polynomial
kernel.  Step 9c freezes the API and safe reuse semantics before the numeric
batch-buffer/GPU refactor.

%===============================================================================
\section{Developer setup and regression commands}
%===============================================================================

From the project root, run the complete suite:
\begin{lstlisting}[language=bash]
python -m unittest discover -s tests -p 'test*.py'
\end{lstlisting}

For the final 0.21.0.dev1 checkpoint the expected gate is:
\begin{lstlisting}
Ran 223 tests
OK
\end{lstlisting}

Run the repeated-query audit:
\begin{lstlisting}[language=bash]
python tools/audit_manybox.py --boxes 40 --workers 4
\end{lstlisting}

Build a wheel without network build isolation when the environment already has
setuptools:
\begin{lstlisting}[language=bash]
python -m pip wheel . --no-deps --no-build-isolation -w dist
\end{lstlisting}

Do not accept a performance change based only on a microbenchmark.  Run the
historical regression suite plus the relevant independent exact/high-precision
audit and an end-to-end engineering benchmark.

%===============================================================================
\section{Source tree for developers}
%===============================================================================

\begin{longtable}{p{0.30\textwidth}p{0.61\textwidth}}
\toprule
Location & What to change there \\
\midrule
\endhead
\codepath{autoia/api.py} & Public one-box/subdivision signatures, validation, exception translation. \\
\codepath{autoia/session.py} & Prepared session, cache, repeated-query orchestration, CPU parallelism. \\
\codepath{autoia/trusted/} & Only proof-critical binary64 primitives and audited reductions/transcendentals.  Keep dependencies minimal. \\
\codepath{autoia/polynomial/} & Sparse polynomial algebra, ranges, Bernstein, recentering. \\
\codepath{autoia/enclosures/} & Representation semantics and arithmetic. \\
\codepath{autoia/semantic/linear_solve.py} & Mathematical linear-system semantic primitive and its residual/contraction proof. \\
\codepath{autoia/semantic/dispatch.py} & Private NumPy/math syntax routing. \\
\codepath{autoia/semantic/scipy_dispatch.py} & Public SciPy syntax compatibility. \\
\codepath{autoia/controller/branching.py} & Uncertain Python paths and piecewise execution. \\
\codepath{autoia/controller/subdivision.py} & Prescribed/adaptive subdivision execution. \\
\codepath{autoia/controller/split_scoring.py} & Cheap non-proof split-direction prediction. \\
\codepath{autoia/elementary/taylor.py} & Taylor coefficients/remainders for smooth functions. \\
\codepath{tests/} & Regression, boundary, exact-rational, high-precision, and integration tests. \\
\codepath{tools/} & Larger randomized/audit/performance scripts that are not shipped as proof dependencies. \\
\bottomrule
\end{longtable}

%===============================================================================
\section{Adding a new supported elementary function}
%===============================================================================

A recommended development sequence is:
\begin{enumerate}[leftmargin=*]
\item State the real mathematical function and domain.
\item Implement a trusted scalar/interval kernel or a trusted coefficient/domain
bound in \codepath{autoia/trusted/elementary.py}.
\item Add independent high-precision point and interval tests.
\item Add Taylor coefficient intervals and a rigorous tail bound in
\codepath{autoia/elementary/taylor.py}.
\item Connect the function to representation-level composition in
\codepath{autoia/enclosures/elementary.py}.
\item Add NumPy/math/SciPy syntax only if it corresponds to a real public API
users commonly write.
\item Add polynomial, interval, generator, object-array, and composition tests.
\item Add a standalone audit if the function is mathematically delicate.
\end{enumerate}

Do not make a platform library's ordinary transcendental return value a trusted
endpoint merely because it is accurate in common cases.

%===============================================================================
\section{Adding a new semantic primitive}
%===============================================================================

For a future operation such as a nonlinear root, eigensystem, or ODE solver:
\begin{enumerate}[leftmargin=*]
\item Define the exact mathematical object being certified.
\item Decide whether the goal is to certify the program algorithm as written or
replace a recognized library call by its mathematical semantics.
\item Build an ordinary floating-point proposal if useful.
\item Derive an independent proof identity, e.g. residual/fixed-point/Krawczyk/
energy argument.
\item Implement the proof in \codepath{autoia/semantic/} using trusted low-level
bounds.
\item Add local dispatch for the recognized library syntax.
\item Preserve a generic validated fallback when reasonable; otherwise fail with
\texttt{UnsupportedOperationError} or \texttt{CertificationError} as appropriate.
\item Add adversarial and near-singular tests, not only nominal examples.
\end{enumerate}

The key rule is: \emph{semantic recognition may change how the operation is
computed, but not what mathematical claim is proved}.

%===============================================================================
\section{Adding or changing a representation}
%===============================================================================

Every new scalar representation must answer four questions:
\begin{enumerate}[leftmargin=*]
\item What exact set \(\Gamma_E(q)\) does the object represent?
\item How is every supported arithmetic operation guaranteed to map input sets to
an output superset?
\item How are conversions to/from existing representations inclusion-preserving?
\item How is a final scalar interval range obtained rigorously?
\end{enumerate}

Only after these are answered should the representation receive a promotion
priority in \texttt{ValidatedValue} dispatch.

For the shared-generator family, never merge two generator IDs solely because
their coefficients or radii are numerically equal.  Equal IDs mean a proved
shared auxiliary variable; different IDs mean independent auxiliary coordinates.

%===============================================================================
\section{Changing trusted arithmetic}
%===============================================================================

Treat \codepath{autoia/trusted/} as a small numerical proof library.

Before changing it:
\begin{itemize}[leftmargin=*]
\item identify the exact-real quantity the function must contain;
\item preserve or extend runtime environment audits;
\item add exact-rational tests for algebraic operations;
\item add independent high-precision tests for transcendental operations;
\item test overflow, subnormal, zero, sign, and endpoint-adversarial cases;
\item keep the historical scalar oracle available when introducing a faster
vectorized path.
\end{itemize}

Do not call higher AutoIA layers from the trusted package.  Architecture boundary
tests are designed to catch such inversions.

%===============================================================================
\section{Changing the many-box cache or parallel backend}
%===============================================================================

The current cache uses a bounded LRU dictionary plus a linear scan for the
smallest containing certificate.  A future spatial index can replace this
without changing the public \texttt{PreparedAutoIA} API.

However, any new cache must preserve these rules:
\begin{enumerate}[leftmargin=*]
\item Cache entries are scoped to one fixed prepared model/configuration.
\item Exact endpoint equality may reuse a certificate directly.
\item Containment-based reuse must verify the new \emph{represented centered
symmetric domain}, not merely the raw user interval, lies in the parent certified
domain.
\item Polynomial recentering must capture coordinate-shift and coefficient
rounding rigorously.
\item Per-request metadata should not mutate result objects previously handed to
users.
\item A cache miss must fall back to a normal fresh proof, not to an approximate
answer.
\end{enumerate}

For process parallelism, remember that every worker owns a separate cache.  A
future shared spatial certificate store will need explicit synchronization and
serialization rules; it should not be introduced implicitly through global
mutable proof state.

%===============================================================================
\section{GPU development guidance}
%===============================================================================

The current user-facing prepared/batch API is intentionally compatible with a
future GPU backend, but the current \texttt{dtype=object} representation is not.
A serious GPU implementation should introduce numeric structure-of-arrays buffers
behind the existing tracing objects.

A sensible staged plan is:
\begin{enumerate}[leftmargin=*]
\item GPU proposal kernels first: batched nominal solves, predictor coefficients,
preconditioner proposals, controller estimates.  Revalidate on the CPU trusted
backend.
\item Numeric CPU batch buffers next, so the same memory model can be transferred
to the GPU.
\item Only later implement a proof-producing GPU backend using explicitly audited
directed-rounding primitives.  Do not assume ordinary CuPy/cuBLAS reductions are
rigorous.
\end{enumerate}

GPU support should plug into \texttt{PreparedAutoIA.certify\_many} or a compatible
backend selection interface rather than introduce an unrelated user API.

%===============================================================================
\section{Release checklist for a code change}
%===============================================================================

Before treating a change as a checkpoint:
\begin{enumerate}[leftmargin=*]
\item Run the complete unit/regression suite.
\item Run the relevant exact-rational or high-precision audit.
\item Run at least one end-to-end engineering smoke test.
\item If performance is claimed, compare against the previous checkpoint on the
same workload and report both kernel and end-to-end timing.
\item Build the wheel from the final source tree.
\item Install that wheel into a clean target and run a smoke test through the
installed package rather than the source checkout.
\item Inspect wheel contents to ensure tests/development scaffolding are not
shipped unintentionally.
\item Update the mathematical contract, code map, user manual, and roadmap when
the public semantics change.
\item Keep a patch archive from the previous checkpoint as well as a complete
source archive.
\end{enumerate}

%===============================================================================
\section{Complete copyable example: nonlinear solve, branch, statistics, and repeated queries}
%===============================================================================

The following example uses several current capabilities together.

\begin{lstlisting}[language=Python]
import numpy as np
from scipy.stats import norm

from autoia import Interval, prepare


def engineering_model(p):
    x, y = p

    A = np.zeros((2, 2))
    A[0, 0] = 2.0 + 0.10*x*x
    A[0, 1] = 0.20
    A[1, 0] = -0.10*y
    A[1, 1] = 1.50 + 0.05*np.exp(x)

    b = np.zeros(2)
    b[0] = 1.0 + 0.10*y
    b[1] = 0.5 + 0.05*np.sin(x)

    u = np.linalg.solve(A, b)

    if x >= 0.0:
        response = u[0] + u[1]*u[1]
    else:
        response = u[0] - u[1]

    probability = norm.cdf(response, loc=0.5, scale=0.2)
    return np.array([response, probability])


session = prepare(
    engineering_model,
    order=4,
    representation="generator",
    cache_max_entries=256,
)

master = session.certify([
    Interval(-0.5, 0.5),
    Interval(-0.2, 0.2),
])

print("master lower:", master.lower)
print("master upper:", master.upper)

queries = [
    [[-0.45, -0.30], [-0.10, 0.00]],
    [[-0.20, -0.05], [ 0.00, 0.10]],
    [[ 0.05,  0.15], [-0.05, 0.05]],
    [[ 0.30,  0.40], [ 0.05, 0.15]],
]

lower, upper = session.certify_many_bounds(
    queries,
    parallel="serial",
)

print("lower bounds:")
print(lower)
print("upper bounds:")
print(upper)
print("session stats:")
print(session.stats)
\end{lstlisting}

For a module-level model and many unrelated boxes, change
\texttt{parallel="serial"} to \texttt{parallel="process"} and specify a worker
count.  If the boxes lie inside a cached parent and are eligible for certificate
restriction, a shared serial session can instead avoid the model/solve execution
entirely for those queries.

%===============================================================================
\section{Final operational guidance}
%===============================================================================

AutoIA is most reliable when the user treats a returned certificate and a failed
certificate differently.  A returned result is a mathematical claim about the
supported program semantics on the declared box.  A failure is not automatically
a statement that the engineering model is invalid; it may mean that the current
representation, domain information, controller budget, or semantic certificate
is insufficiently sharp.

For development, preserve the same distinction in code: proposal logic may be
optimized aggressively, but every accepted result must remain traceable to the
trusted enclosure and semantic proof contracts.  The prepared-session layer adds
a second important rule: reuse is safe only when a previously proved certificate
is applied to a domain that has itself been rigorously shown to lie inside the
parent certified domain.

\end{document}
